\documentclass[11pt]{article}
\usepackage[margin=1in]{geometry}
\usepackage{amsmath, amssymb, bm}
\usepackage{booktabs}

\title{Cost-Push Markup Shocks in the Rubbo Network Phillips Curve}
\author{}
\date{}

\begin{document}
\maketitle

\section{Baseline Objects}

Let there be $N$ sectors. The calibrated primitives are
\[
    \alpha \in \mathbb{R}^{N}, \qquad
    \beta \in \mathbb{R}^{N}, \qquad
    \Omega \in \mathbb{R}^{N \times N}, \qquad
    \Delta = \mathrm{diag}(\delta_i).
\]
Here $\alpha_i$ is the labor share in sector $i$, $\beta_i$ is the final-consumption
weight, $\Omega_{ij}$ is the input share of good $j$ in sector $i$, and $\delta_i$
is the Calvo price-adjustment probability.

The Domar/sales exposure vector is
\[
    \lambda^\prime
    =
    \beta^\prime (I-\Omega)^{-1}.
\]
Define the rigidity-adjusted Leontief inverse
\[
    R_\Delta = (I-\Omega\Delta)^{-1},
\]
and the common denominator
\[
    d
    =
    1 - \beta^\prime \Delta R_\Delta \alpha.
\]

The sectoral Phillips-curve slope vector, following Rubbo's Matlab convention,
is
\[
    b
    =
    (\gamma+\varphi)
    \frac{\Delta R_\Delta \alpha}{d}.
\]
The CPI Phillips-curve slope is
\[
    \kappa^C = \beta^\prime b.
\]

\section{Productivity-Shock Term in the Baseline Model}

For reference, the productivity-shock propagation matrix used in the original code is
\[
    v^A
    =
    \Delta R_\Delta
    \left[
        \alpha
        \frac{\lambda^\prime - \beta^\prime \Delta R_\Delta}{d}
        - I
    \right].
\]
The CPI residual generated by productivity shocks is
\[
    u^C_t = \beta^\prime v^A \, d\log A_t.
\]

\section{Markup Cost-Push Shocks}

Following the supplement, introduce an exogenous desired-markup shock
\[
    \mu_t^D \in \mathbb{R}^{N}.
\]
A positive element of $\mu_t^D$ is a cost-push shock in that sector. A sectoral
subsidy can be represented as a negative markup shock:
\[
    \mu^{D,\mathrm{net}}_t
    =
    \mu^{D,\mathrm{shock}}_t
    -
    s_t.
\]

The consumer-price markup term is
\[
    v_t^C
    =
    \frac{
        \beta^\prime \Delta (I-\Omega\Delta)^{-1}
    }{
        1-\beta^\prime \Delta (I-\Omega\Delta)^{-1}\alpha
    }
    \mu_t^D.
\]
This is the object computed in the static module as \texttt{CPI\_markup}.

The divine-coincidence markup term is
\[
    \pi_t^{DC}
    =
    \rho E_t \pi_{t+1}^{DC}
    +(\gamma+\varphi)\tilde y_t
    + \lambda^\prime \mu_t^D.
\]
Therefore, absent expectations and with $\tilde y_t=0$, the static DC response is
\[
    \lambda^\prime \mu_t^D.
\]
The output gap that stabilizes the divine-coincidence index on impact is
\[
    \tilde y_t^{DC=0}
    =
    -\frac{\lambda^\prime \mu_t^D}{\gamma+\varphi}.
\]

\section{Sectoral Phillips Curves}

Let
\[
    B =
    \frac{\Delta R_\Delta \alpha}{d}
\]
so that $b=(\gamma+\varphi)B$. Rubbo's sectoral matrix is
\[
V
=
    \left[
        \Delta R_\Delta
        -
        B
        \left(
            \lambda^\prime
            -
            \beta^\prime \Delta R_\Delta
        \right)
    \right]
    (I-\Omega).
\]
The sectoral Phillips curves with desired-markup shocks are
\[
    \pi_t
    =
    \rho (I-V)E_t\pi_{t+1}
    +
    b \tilde y_t
    -
    V\chi_t
    +
    (I-V)\Delta(I-\Delta)^{-1}\mu_t^D.
\]
In the current markup-shock module, the productivity term is set to zero,
$\chi_t=0$, so the dynamic recursion is
\[
    \pi_t
    =
    \rho (I-V)E_t\pi_{t+1}
    +
    b \tilde y_t
    +
    (I-V)\Delta(I-\Delta)^{-1}\mu_t^D.
\]

\section{Dynamic Experiment}

The current dynamic script assumes an exogenous AR(1) process:
\[
    \mu_{t+1}^D = \rho_\mu \mu_t^D.
\]
Given a terminal condition $\pi_{T+1}=0$, the sectoral Phillips curve is solved
backwards:
\[
    \pi_t
    =
    \rho (I-V)\pi_{t+1}
    +
    b\tilde y_t
    +
    (I-V)\Delta(I-\Delta)^{-1}\mu_t^D.
\]
The reported CPI inflation response is
\[
    \pi_t^C = \beta^\prime \pi_t,
\]
and the DC response is computed from the scalar recursion
\[
    \pi_t^{DC}
    =
    \rho \pi_{t+1}^{DC}
    +
    (\gamma+\varphi)\tilde y_t
    +
    \lambda^\prime \mu_t^D.
\]

\section{Interest-Rate Implementation}

The output-gap rules above can be translated into nominal interest-rate paths using
Rubbo's log-linear IS curve. With consumer-price inflation entering the household Euler
equation, the IS curve is
\[
    \tilde y_t
    =
    E_t \tilde y_{t+1}
    -
    \frac{1}{\gamma}
    \left(
        i_{t+1}
        -
        E_t \pi^C_{t+1}
        -
        r^{nat}_{t+1}
    \right).
\]
Rearranging gives the policy-rate response relative to the natural nominal rate:
\[
    i_{t+1}-r^{nat}_{t+1}
    =
    E_t\pi^C_{t+1}
    +
    \gamma
    \left(
        E_t\tilde y_{t+1}-\tilde y_t
    \right).
\]
The code uses the simulated future paths as the model-consistent expectations:
\[
    E_t\pi^C_{t+1} = \pi^C_{t+1},
    \qquad
    E_t\tilde y_{t+1} = \tilde y_{t+1},
\]
with terminal values set to zero at the end of the finite horizon. Therefore the
computed interest-rate path is
\[
    \widehat i_{t+1}
    \equiv
    i_{t+1}-r^{nat}_{t+1}
    =
    \pi^C_{t+1}
    +
    \gamma
    \left(
        \tilde y_{t+1}-\tilde y_t
    \right).
\]
Because the oil experiment is a desired-markup shock rather than a productivity shock,
the productivity component of $r^{nat}_{t+1}$ is unchanged. Thus $\widehat i_{t+1}$
is the nominal policy-rate movement above the natural-rate path.
When reporting peak interest-rate movements, the code excludes the final finite-horizon
transition by default, because the last point may reflect the imposed terminal condition
rather than the economic response. The full path is still stored.

\section{Policy Rules Implemented}

The first dynamic module implements two simple output-gap rules:
\[
    \tilde y_t = 0
    \qquad \text{or} \qquad
    \tilde y_t = -\frac{\lambda^\prime \mu_t^D}{\gamma+\varphi}.
\]
The first isolates pass-through and inflation dynamics. The second stabilizes the
divine-coincidence Phillips curve period by period. A full Taylor-rule block can be
added later by combining the markup-shock recursion with Rubbo's monetary IRF
system.

\section{Policy Comparison Module}

The policy-comparison script expresses monetary policy as a path for the output gap.
This is a reduced-form policy experiment: it asks what output-gap movement would be
needed under alternative targets, before translating that movement into an interest-rate
rule through an IS curve.

The rules currently implemented are:
\[
    \tilde y_t^{0}=0,
\]
\[
    \tilde y_t^{DC}
    =
    -\frac{\lambda^\prime \mu_t^D}{\gamma+\varphi},
\]
and a CPI-stabilizing rule. Let $\pi_{t+1}$ denote the sectoral inflation vector already
computed in the backward recursion. The CPI-stabilizing output gap is chosen so that
$\beta^\prime \pi_t=0$:
\[
    \tilde y_t^{CPI}
    =
    -
    \frac{
        \rho \beta^\prime (I-V)\pi_{t+1}
        +
        \beta^\prime (I-V)\Delta(I-\Delta)^{-1}\mu_t^D
    }{
        \beta^\prime b
    }.
\]
The script also includes partial-reaction rules:
\[
    \tilde y_t^{lean,CPI}
    =
    \theta_C \tilde y_t^{CPI},
    \qquad
    \tilde y_t^{lean,DC}
    =
    \theta_{DC} \tilde y_t^{DC},
\]
where $\theta_C,\theta_{DC}\in[0,1]$.

The reported losses are transparent quadratic losses, not the full welfare matrix from
Rubbo's Section 5:
\[
    \mathcal{L}^{CPI}
    =
    \sum_{t=0}^{T-1}
    \rho_L^t
    \left[
        (\pi_t^C)^2
        +
        \omega_y \tilde y_t^2
    \right],
\]
\[
    \mathcal{L}^{DC}
    =
    \sum_{t=0}^{T-1}
    \rho_L^t
    \left[
        (\pi_t^{DC})^2
        +
        \omega_y \tilde y_t^2
    \right],
\]
and
\[
    \mathcal{L}^{dual}
    =
    \sum_{t=0}^{T-1}
    \rho_L^t
    \left[
        \omega_C(\pi_t^C)^2
        +
        \omega_{DC}(\pi_t^{DC})^2
        +
        \omega_y \tilde y_t^2
    \right].
\]
These losses are intended to compare simple central-bank responses to a markup shock:
ignore the shock, stabilize CPI, stabilize the divine-coincidence index, or lean
partially against either target.

\section{Code Procedures}

The project is organized in layers. The prototype files are left unchanged. The new
markup-shock layer adds:
\begin{enumerate}
    \item \texttt{markup\_costpush\_static.m}: computes static pass-through objects,
    including CPI markup response, divine-coincidence response, the output gap that
    stabilizes divine-coincidence inflation, network pressure, and the sectoral
    Phillips-curve markup term.
    \item \texttt{markup\_costpush\_irf.m}: builds an AR(1) path for $\mu_t^D$ and
    solves the forward-looking sectoral Phillips curves backward over a finite horizon.
    \item \texttt{markup\_costpush\_policy\_compare.m}: compares simple output-gap
    policy rules and computes transparent quadratic losses.
    \item \texttt{markup\_policy\_rate\_path.m}: converts each output-gap path into a
    nominal interest-rate path relative to the natural rate using the IS curve above.
\end{enumerate}
The main drivers are:
\begin{enumerate}
    \item \texttt{run\_oil\_markup\_module\_example.m}: static and dynamic pass-through
    under a baseline oil markup shock and an illustrative subsidy.
    \item \texttt{oil\_markup\_monetary\_policy.m}: compares output-gap policy rules,
    inflation responses, output-gap paths, quadratic losses, and interest-rate paths.
    \item \texttt{oil\_markup\_interest\_rate\_policy.m}: focuses on the interest-rate
    implementation of the same policy rules.
\end{enumerate}
All drivers follow Rubbo-style paths and can be run from the replication root, from
\texttt{calibration/oil shocks}, or from \texttt{calibration/data cleaning/IO tables}.

\end{document}
